% 康奈尔听课笔记模板（配色取自王道课件）
\documentclass[10.5pt,a4paper,oneside]{article}
\usepackage[a4paper,top=2.3cm,bottom=2.2cm,left=5.7cm,right=1.5cm,
  marginparwidth=4.25cm,marginparsep=0.55cm,headheight=14pt,footskip=1.1cm]{geometry}
\usepackage{xeCJK}
\usepackage{xeCJKfntef}
\xeCJKDeclareCharClass{FullLeft}{"201C,"2018}
\xeCJKDeclareCharClass{FullRight}{"201D,"2019}
\xeCJKDeclareCharClass{CJK}{"2460->"24FF,"2605,"2606}
\setmainfont{Noto Sans CJK SC}[BoldFont={Noto Sans CJK SC Bold},ItalicFont={Noto Sans CJK SC},ItalicFeatures={FakeSlant=0.18}]
\setsansfont{Noto Sans CJK SC}[BoldFont={Noto Sans CJK SC Bold}]
\setmonofont{DejaVu Sans Mono}[Scale=0.9]
\setCJKmainfont{Noto Sans CJK SC}[BoldFont={Noto Sans CJK SC Bold},ItalicFont={Noto Sans CJK SC},ItalicFeatures={FakeSlant=0.18}]
\setCJKsansfont{Noto Sans CJK SC}[BoldFont={Noto Sans CJK SC Bold}]
\setCJKmonofont{Noto Sans Mono CJK SC}[Scale=0.9]
\linespread{1.32}
\usepackage{amsmath,amssymb}
\usepackage[dvipsnames,table]{xcolor}
\usepackage{graphicx}
\usepackage{adjustbox}
\usepackage{tikz}
\usetikzlibrary{calc}
\usepackage[most]{tcolorbox}
\usepackage{eso-pic}
\usepackage{fancyhdr}
\usepackage{titlesec}
\usepackage{needspace}
\usepackage{enumitem}
\usepackage{xltabular,booktabs,array,colortbl}
\usepackage{marginnote}
\reversemarginpar
\renewcommand*{\raggedleftmarginnote}{}
\renewcommand*{\raggedrightmarginnote}{}
\usepackage{caption}
\usepackage{titletoc}
\usepackage{fancyvrb}
\usepackage{framed}
\usepackage[section]{placeins}
\usepackage{hyperref}
\hypersetup{colorlinks=true,linkcolor=wdblue,urlcolor=wdblue,pdftitle={2.3.5
静态链表},pdfauthor={听课笔记}}

% ---------- 课件配色 ----------
\definecolor{wdblue}{HTML}{4472C4}
\definecolor{wdlightblue}{HTML}{5B9BD5}
\definecolor{wdtitle}{HTML}{4A86C8}
\definecolor{wddark}{HTML}{1F3864}
\definecolor{wdorange}{HTML}{ED7D31}
\definecolor{wdgray}{HTML}{A5A5A5}
\definecolor{wdgold}{HTML}{FFC000}
\definecolor{wdgreen}{HTML}{70AD47}
\definecolor{wddeepgreen}{HTML}{3F7F2A}
\definecolor{wdred}{HTML}{C00000}
\definecolor{wdpurple}{HTML}{7030A0}
\definecolor{cuebg}{HTML}{F2F6FC}
\definecolor{codebg}{HTML}{F7F9FC}
\color[HTML]{262626}

% ---------- 行内样式 ----------
\newcommand{\hlred}[1]{\textcolor{wdred}{\bfseries #1}}
\newcommand{\hlblue}[1]{\textcolor{wdblue}{\bfseries #1}}
\newcommand{\hlgreen}[1]{\textcolor{wddeepgreen}{\bfseries #1}}
\newcommand{\hlorange}[1]{\textcolor{wdorange}{\bfseries #1}}
\newcommand{\hlpurple}[1]{\textcolor{wdpurple}{\bfseries #1}}
\newcommand{\uln}[1]{\CJKunderline[format=\color{wdred},thickness=0.9pt,depth=0.28em]{#1}}
\newcommand{\wavy}[1]{\CJKunderwave[format=\color{wdorange}]{#1}}
\newcommand{\hlmark}[1]{\CJKunderline[format=\color{wdgold!55},thickness=0.55em,depth=-0.05em]{#1}}
\newcommand{\tagbox}[2]{\tcbox[on line,colback=#1,colframe=#1,coltext=white,boxsep=0.5pt,left=3pt,right=3pt,top=0.5pt,bottom=0.5pt,arc=2pt,fontupper=\small\bfseries]{#2}}

% ---------- 康奈尔版式背景：左侧线索栏 ----------
\newlength{\cuestripw}\setlength{\cuestripw}{5.15cm}
\AddToShipoutPictureBG{%
  \begin{tikzpicture}[remember picture,overlay]
    \fill[cuebg] (current page.north west) rectangle ($(current page.south west)+(\cuestripw,0)$);
    \draw[wdblue,line width=1.2pt] ($(current page.north west)+(\cuestripw,-1.55cm)$) -- ($(current page.south west)+(\cuestripw,1.55cm)$);
    \node[anchor=west,text=wdblue,font=\small\bfseries] at ($(current page.north west)+(0.75cm,-1.25cm)$) {线索 · 问题};
    \node[anchor=west,text=wdblue,font=\small\bfseries] at ($(current page.north west)+(\cuestripw+0.55cm,-1.25cm)$) {笔记 · 老师主线 + 补充讲解};
  \end{tikzpicture}}

% ---------- 页眉页脚 ----------
\pagestyle{fancy}\fancyhf{}
\renewcommand{\headrulewidth}{0pt}
\fancyhead[R]{\small\color{wdgray}第2章 线性表}
\fancyfoot[C]{\small\color{wdgray}\thepage}
\fancyfoot[R]{\small\color{wdorange}2.3.5 静态链表 · 听课笔记}
\fancyfoot[L]{\small\hyperlink{toc}{\textcolor{wdblue}{\ensuremath{\uparrow}\,回到目录}}}
\fancypagestyle{plain}{\fancyhf{}\fancyfoot[C]{\small\color{wdgray}\thepage}\fancyfoot[R]{\small\color{wdorange}2.3.5
静态链表 · 听课笔记}}

% ---------- 标题：仿课件的蓝色标题 + 六色条 ----------
\newcommand{\sixbar}{\par\nointerlineskip\vspace{3pt}\centering
  \begin{tikzpicture}\foreach \c [count=\i] in {wdlightblue,wdorange,wdgray,wdgold,wdblue,wdgreen}
     {\fill[\c] ({(\i-1)*1.05},0) rectangle ({(\i-1)*1.05+0.92},0.11);}\end{tikzpicture}\par}
\titleformat{\section}[block]{\centering\Large\bfseries\color{wdtitle}}{}{0pt}{}[\sixbar]
\titlespacing*{\section}{0pt}{16pt}{10pt}
\pretocmd{\section}{\Needspace{8\baselineskip}}{}{}
\titleformat{\subsection}{\Needspace{7\baselineskip}\large\bfseries\color{wddark}}{\textcolor{wdblue}{\rule[-0.15em]{4pt}{1.05em}}\hspace{6pt}}{0pt}{}
\titlespacing*{\subsection}{0pt}{12pt}{5pt}
\titleformat{\subsubsection}{\Needspace{5\baselineskip}\normalsize\bfseries\color{wdblue}}{\textcolor{wdorange}{\large\ensuremath{\blacktriangleright}}\hspace{4pt}}{0pt}{}
\titlespacing*{\subsubsection}{0pt}{8pt}{3pt}
\setcounter{secnumdepth}{0}
\setlength{\parindent}{0pt}\setlength{\parskip}{5pt}
\setlist{nosep,leftmargin=1.5em,itemsep=2pt}
\setlist[itemize,1]{label=\textcolor{wdblue}{\small\ensuremath{\bullet}}}
\setlist[itemize,2]{label=\textcolor{wdorange}{\scriptsize\ensuremath{\blacktriangleright}}}
\setlist[enumerate,1]{label=\textcolor{wdblue}{\bfseries\arabic*.}}
\providecommand{\tightlist}{\setlength{\itemsep}{2pt}\setlength{\parskip}{0pt}}

% ---------- 目录 ----------
\setcounter{tocdepth}{1}
\titlecontents{section}[0pt]{\addvspace{5pt}\bfseries\color{wddark}}{}{}{\titlerule*[0.7pc]{\color{wdgray!60}.}\contentspage}
\titlecontents{subsection}[1.6em]{\addvspace{1pt}\small\color{black!75}}{}{}{\titlerule*[0.7pc]{\color{wdgray!40}.}\contentspage}
\newcommand{\tocback}{\par\vspace{-3pt}{\hfill\footnotesize\hyperlink{toc}{\textcolor{wdgray}{\ensuremath{\uparrow}\,回到目录}}}\par\vspace{-4pt}}

% ---------- 线索栏 ----------
\newcommand{\cueblock}[2]{\Needspace{4\baselineskip}\leavevmode#1{#2}\par\vspace{-\baselineskip}\vspace{-\parskip}}
% 线索直接贴在段首行的左侧（不用 marginnote，避免分页时偶尔落进正文栏）
\newcommand{\cuebox}[1]{\leavevmode\llap{\smash{\begin{minipage}[t]{\marginparwidth}\raggedright\small\linespread{1.18}\selectfont #1\end{minipage}}\hspace{\marginparsep}}}
\newcommand{\cuenote}[1]{\cuebox{{\color{wdblue}\bfseries Q\,}\color{wddark}#1}}
\newcommand{\cuekey}[1]{\cuebox{\color{wdred}\bfseries #1}}

% ---------- 讲解框 ----------
\tcbset{wdbox/.style={enhanced,breakable,arc=2pt,boxrule=0pt,leftrule=3pt,left=7pt,right=7pt,top=5pt,bottom=5pt,
   before skip=7pt,after skip=7pt,fonttitle=\bfseries\small,attach title to upper={\quad},coltitle=#1,
   colframe=#1,colback=#1!6}}
\newtcolorbox{defbox}[1][]{enhanced,breakable,arc=3pt,boxrule=0.8pt,colframe=wdblue,colback=wdblue!4,
   left=8pt,right=8pt,top=6pt,bottom=6pt,fonttitle=\bfseries\small,coltitle=white,colbacktitle=wdblue,
   title={定义#1},before skip=8pt,after skip=8pt}
\newtcolorbox{bridgebox}{wdbox=wdblue,title={承上启下}}
\newtcolorbox{plainbox}{wdbox=wddeepgreen,colframe=wdgreen,colback=wdgreen!8,title={大白话}}
\newtcolorbox{deepbox}{wdbox=wdpurple,title={背后的思想}}
\newtcolorbox{exambox}{wdbox=wdred,title={考点}}
\newtcolorbox{pitfallbox}{wdbox=wdorange,title={易错提醒}}
\newtcolorbox{bookbox}{wdbox=wddark,colback=wdgray!10,colframe=wdgray,title={王道书补充}}
\newtcolorbox{teacherbox}{wdbox=wdorange,colback=wdgold!8,colframe=wdgold,coltitle=wdorange,title={老师原话}}
\newtcolorbox{fixbox}{wdbox=wdred,colback=wdred!4,title={勘误 / 说明}}
% 习题笔记用
\newtcolorbox{questionbox}[1][]{enhanced,breakable,arc=3pt,boxrule=0.8pt,colframe=wddark,colback=white,
   left=8pt,right=8pt,top=6pt,bottom=6pt,fonttitle=\bfseries\small,coltitle=white,colbacktitle=wddark,
   title={#1},before skip=10pt,after skip=6pt}
\newtcolorbox{videobox}{wdbox=wdorange,colback=wdgold!8,colframe=wdgold,coltitle=wdorange,title={视频讲解}}
\newtcolorbox{jiexibox}{wdbox=wddark,colback=wdgray!10,colframe=wdgray,title={书上解析}}
\newtcolorbox{addbox}{wdbox=wdpurple,title={我的补充}}
\newtcolorbox{recallbox}{wdbox=wdblue,title={回顾正课}}
\newtcolorbox{watchbox}{enhanced,breakable,arc=3pt,boxrule=1pt,colframe=wdorange,colback=wdgold!12,
   fonttitle=\bfseries\small,coltitle=white,colbacktitle=wdorange,title={建议去看视频},left=8pt,right=8pt,before skip=8pt,after skip=8pt}
\newtcolorbox{summarybox}{enhanced,arc=3pt,
   boxrule=1pt,colframe=wdblue,colback=white,colbacktitle=wdblue,coltitle=white,fonttitle=\bfseries,
   title={总结区 · 用自己的话复述这一节},left=10pt,right=10pt,top=8pt,bottom=8pt,before skip=14pt}

% ---------- 代码 ----------
\usepackage{color}
\usepackage{fancyvrb}
\newcommand{\VerbBar}{|}
\newcommand{\VERB}{\Verb[commandchars=\\\{\}]}
\DefineVerbatimEnvironment{Highlighting}{Verbatim}{commandchars=\\\{\}}
% Add ',fontsize=\small' for more characters per line
\newenvironment{Shaded}{}{}
\newcommand{\AlertTok}[1]{\textcolor[rgb]{1.00,0.00,0.00}{#1}}
\newcommand{\AnnotationTok}[1]{\textcolor[rgb]{0.38,0.63,0.69}{#1}}
\newcommand{\AttributeTok}[1]{\textcolor[rgb]{0.49,0.56,0.16}{#1}}
\newcommand{\BaseNTok}[1]{\textcolor[rgb]{0.77,0.35,0.07}{#1}}
\newcommand{\BuiltInTok}[1]{\textcolor[rgb]{0.00,0.50,0.00}{#1}}
\newcommand{\CharTok}[1]{\textcolor[rgb]{0.63,0.19,0.43}{#1}}
\newcommand{\CommentTok}[1]{\textcolor[rgb]{0.18,0.55,0.23}{#1}}
\newcommand{\CommentVarTok}[1]{\textcolor[rgb]{0.18,0.55,0.23}{#1}}
\newcommand{\ConstantTok}[1]{\textcolor[rgb]{0.77,0.35,0.07}{#1}}
\newcommand{\ControlFlowTok}[1]{\textcolor[rgb]{0.12,0.31,0.71}{#1}}
\newcommand{\DataTypeTok}[1]{\textcolor[rgb]{0.12,0.31,0.71}{#1}}
\newcommand{\DecValTok}[1]{\textcolor[rgb]{0.77,0.35,0.07}{#1}}
\newcommand{\DocumentationTok}[1]{\textcolor[rgb]{0.18,0.55,0.23}{#1}}
\newcommand{\ErrorTok}[1]{\textcolor[rgb]{1.00,0.00,0.00}{#1}}
\newcommand{\ExtensionTok}[1]{#1}
\newcommand{\FloatTok}[1]{\textcolor[rgb]{0.77,0.35,0.07}{#1}}
\newcommand{\FunctionTok}[1]{\textcolor[rgb]{0.44,0.19,0.63}{#1}}
\newcommand{\ImportTok}[1]{\textcolor[rgb]{0.44,0.19,0.63}{#1}}
\newcommand{\InformationTok}[1]{\textcolor[rgb]{0.38,0.63,0.69}{#1}}
\newcommand{\KeywordTok}[1]{\textcolor[rgb]{0.12,0.31,0.71}{#1}}
\newcommand{\NormalTok}[1]{#1}
\newcommand{\OperatorTok}[1]{\textcolor[rgb]{0.20,0.20,0.20}{#1}}
\newcommand{\OtherTok}[1]{\textcolor[rgb]{0.00,0.44,0.13}{#1}}
\newcommand{\PreprocessorTok}[1]{\textcolor[rgb]{0.44,0.19,0.63}{#1}}
\newcommand{\RegionMarkerTok}[1]{#1}
\newcommand{\SpecialCharTok}[1]{\textcolor[rgb]{0.63,0.19,0.43}{#1}}
\newcommand{\SpecialStringTok}[1]{\textcolor[rgb]{0.63,0.19,0.43}{#1}}
\newcommand{\StringTok}[1]{\textcolor[rgb]{0.63,0.19,0.43}{#1}}
\newcommand{\VariableTok}[1]{\textcolor[rgb]{0.10,0.09,0.49}{#1}}
\newcommand{\VerbatimStringTok}[1]{\textcolor[rgb]{0.63,0.19,0.43}{#1}}
\newcommand{\WarningTok}[1]{\textcolor[rgb]{0.38,0.63,0.69}{#1}}
\makeatletter
\@ifundefined{Shaded}{}{\renewenvironment{Shaded}{\def\FrameCommand{\fboxsep=5pt\fboxrule=0.5pt\fcolorbox{wdblue!35}{codebg}}%
   \MakeFramed{\advance\hsize-\width\FrameRestore}\small}{\endMakeFramed}}
\makeatother
\fvset{fontsize=\small}

% ---------- 表格 ----------
\renewcommand{\arraystretch}{1.42}
\newcolumntype{Y}{>{\raggedright\arraybackslash}X}

% ---------- 图片 ----------
% 课件截图做成浮动体：LaTeX 会自动用正文填满页面，避免大块空白
\setcounter{topnumber}{3}\setcounter{bottomnumber}{2}\setcounter{totalnumber}{4}
\renewcommand{\topfraction}{0.85}\renewcommand{\bottomfraction}{0.6}\renewcommand{\textfraction}{0.12}\renewcommand{\floatpagefraction}{0.75}
\setlength{\textfloatsep}{10pt plus 3pt minus 2pt}\setlength{\floatsep}{8pt plus 2pt}
\newcommand{\slidefig}[3]{\begin{figure}[!ht]\centering\begin{tcolorbox}[enhanced,width=#2\linewidth,arc=2pt,boxrule=0.6pt,colframe=wdgray!70,
   colback=white,left=2pt,right=2pt,top=2pt,bottom=2pt,drop fuzzy shadow=black!25]\includegraphics[width=\linewidth]{#1}\end{tcolorbox}
   \vspace{-3pt}{\small\color{wdgray}#3}\end{figure}}
\newcommand{\plainfig}[3]{\par\noindent\begin{minipage}{\linewidth}\centering\includegraphics[width=#2\linewidth]{#1}\par\vspace{2pt}{\small\color{wdgray}#3}\end{minipage}\par\vspace{4pt}}

\begin{document}
% ---------- 封面条（仿课件封面） ----------
\thispagestyle{plain}
\begin{tcolorbox}[enhanced,enlarge left by=-4.8cm,width=\linewidth+4.8cm,arc=0pt,boxrule=0pt,
   colback=wddark,interior style={left color=wddark,right color=wdblue},left=18pt,right=18pt,top=14pt,bottom=12pt,
   overlay={\draw[white,line width=2.5pt] ([xshift=10pt,yshift=-10pt]frame.north west) -- ++(0,-0.9cm);
            \draw[white,line width=2.5pt] ([xshift=10pt,yshift=-10pt]frame.north west) -- ++(1.6cm,0);
            \draw[white,line width=2.5pt] ([xshift=-10pt,yshift=10pt]frame.south east) -- ++(0,0.9cm);
            \draw[white,line width=2.5pt] ([xshift=-10pt,yshift=10pt]frame.south east) -- ++(-1.6cm,0);}]
  {\color{white!85!wdlightblue}\small 数据结构 · 第2章
线性表 · 康奈尔听课笔记}\par\vspace{6pt}
  {\color{white}\fontsize{22}{28}\selectfont\bfseries 2.3.5
静态链表}\par\vspace{6pt}
  {\color{white!80!wdlightblue}\small 课件 2.3.5 · 正课字幕 P18 · 王道书
§2.3.5（第 38 页）}
\end{tcolorbox}
\vspace{4pt}
\begin{tcolorbox}[enhanced,arc=3pt,boxrule=0.8pt,colframe=wdblue,colback=wdblue!3,left=10pt,right=10pt,top=6pt,bottom=6pt,
   colbacktitle=wdblue,coltitle=white,fonttitle=\bfseries\small,title={目录 · 点击标题跳转}]
\hypertarget{toc}{}\makeatletter\@starttoc{toc}\makeatother
\end{tcolorbox}
\hypertarget{ux8fd9ux4e00ux8282ux5728ux8bb2ux4ec0ux4e48}{%
\section{这一节在讲什么}\label{ux8fd9ux4e00ux8282ux5728ux8bb2ux4ec0ux4e48}}

\tocback

\Needspace{4\baselineskip}

\cuenote{这一节讲静态链表的哪三件事？}三件事：什么是静态链表；怎么用代码定义一个静态链表；简述基本操作怎么实现。老师说这个考点考得不多，也很少考代码实现，主要想借它补充一些
C 语言的知识，\uln{至少要能看懂书上给的代码}。

\hypertarget{ux4e00ux4ec0ux4e48ux662fux9759ux6001ux94feux8868}{%
\section{一、什么是静态链表}\label{ux4e00ux4ec0ux4e48ux662fux9759ux6001ux94feux8868}}

\tocback

\cueblock{\cuenote}{\relax{静态链表的“指针”是什么？表尾怎么表示？第 \(i\) 号结点的地址怎么算？}}

\begin{itemize}
\tightlist
\item
  \textbf{单链表}：各个结点离散地分布在内存的各个角落（课件：“星罗棋布、散落天涯”），每个结点存数据元素和指向下一个结点的指针（下一个结点的内存地址）。
\item
  \textbf{静态链表}：分配\uln{一整片连续的内存空间}，各个结点集中安置在这片空间中的某些位置。每个结点存数据元素和\textbf{下一个结点的数组下标}。
\end{itemize}

看课件的例子（数组下标 0～9）：

\begin{itemize}
\tightlist
\item
  \uln{0 号结点充当“头结点”}，不存实际的数据元素；它的游标是
  2，所以第一个数据结点（位序为 1）存在下标 2 处（\(e_1\)），\(e_1\)
  的游标是 1，下一个就是下标 1 处的
  \(e_2\)，依此类推：\(e_1 \to e_2 \to e_3 \to e_4\)，存放的下标依次为
  2、1、6、3。
\item
  数组下标也叫\textbf{游标}，作用和单链表的指针差不多：指针指明具体的内存地址，游标只指明下一个元素的数组下标。
\item
  单链表表尾的 \texttt{next} 指向
  \texttt{NULL}；静态链表中\uln{游标为 \(-1\) 表示已经到达表尾}，后面没有结点了。
\end{itemize}

\slidefig{assets/2.3.5-p03.jpg}{0.65}{课件：什么是静态链表}

\textbf{地址怎么算}：各结点的空间连续。设每个数据元素 4B、每个游标
4B（每个结点共 8B），起始地址为 \texttt{addr}，则下标为 2 的结点 \(e_1\)
的存放地址为 \(\texttt{addr}+8\times2\)。即：

\[\text{下标为 } i \text{ 的结点的地址} = \texttt{addr} + \text{每个结点的大小} \times i\]

用这种方式就能把游标（数组下标）映射成对应结点的实际存放地址，这就是静态链表的基本原理。

\hypertarget{ux4e8cux7528ux4ee3ux7801ux5b9aux4e49ux9759ux6001ux94feux8868}{%
\section{二、用代码定义静态链表}\label{ux4e8cux7528ux4ee3ux7801ux5b9aux4e49ux9759ux6001ux94feux8868}}

\tocback

\cueblock{\cuenote}{\relax{课本的 \texttt{typedef\ struct\ \{…\}\ SLinkList{[}MaxSize{]};}
是什么意思？为什么不直接写 \texttt{struct\ Node\ a{[}MaxSize{]}}？}}

\hypertarget{ux5bb9ux6613ux60f3ux5230ux7684ux5199ux6cd5}{%
\subsection{容易想到的写法}\label{ux5bb9ux6613ux60f3ux5230ux7684ux5199ux6cd5}}

定义一个结点结构体，再声明一个结点数组：

\begin{Shaded}
\begin{Highlighting}[]
\PreprocessorTok{\#define MaxSize }\DecValTok{10}\PreprocessorTok{           }\CommentTok{// 静态链表的最大长度}
\KeywordTok{struct}\NormalTok{ Node }\OperatorTok{\{}                \CommentTok{// 静态链表结构类型的定义}
\NormalTok{    ElemType data}\OperatorTok{;}           \CommentTok{// 存储数据元素}
    \DataTypeTok{int}\NormalTok{ next}\OperatorTok{;}                \CommentTok{// 下一个元素的数组下标}
\OperatorTok{\};}

\DataTypeTok{void}\NormalTok{ testSLinkList}\OperatorTok{()} \OperatorTok{\{}
    \KeywordTok{struct}\NormalTok{ Node a}\OperatorTok{[}\NormalTok{MaxSize}\OperatorTok{];}  \CommentTok{// 数组 a 作为静态链表}
    \CommentTok{// ……后续代码}
\OperatorTok{\}}
\end{Highlighting}
\end{Shaded}

\texttt{a} 有 \texttt{MaxSize}（10）个元素，每个元素是一个
\texttt{struct\ Node}，占内存中一整片连续空间。

\hypertarget{ux8bfeux672cux7684ux5199ux6cd5}{%
\subsection{课本的写法}\label{ux8bfeux672cux7684ux5199ux6cd5}}

\begin{Shaded}
\begin{Highlighting}[]
\PreprocessorTok{\#define MaxSize }\DecValTok{10}\PreprocessorTok{           }\CommentTok{// 静态链表的最大长度}
\KeywordTok{typedef} \KeywordTok{struct} \OperatorTok{\{}             \CommentTok{// 静态链表结构类型的定义}
\NormalTok{    ElemType data}\OperatorTok{;}           \CommentTok{// 存储数据元素}
    \DataTypeTok{int}\NormalTok{ next}\OperatorTok{;}                \CommentTok{// 下一个元素的数组下标}
\OperatorTok{\}}\NormalTok{ SLinkList}\OperatorTok{[}\NormalTok{MaxSize}\OperatorTok{];}

\DataTypeTok{void}\NormalTok{ testSLinkList}\OperatorTok{()} \OperatorTok{\{}
\NormalTok{    SLinkList a}\OperatorTok{;}
    \CommentTok{// ……后续代码}
\OperatorTok{\}}
\end{Highlighting}
\end{Shaded}

老师说这种写法挺少见，问了身边好几个常用 C/C++
的同学也很少有人见过。它等价于：

\begin{Shaded}
\begin{Highlighting}[]
\KeywordTok{struct}\NormalTok{ Node }\OperatorTok{\{}
\NormalTok{    ElemType data}\OperatorTok{;}
    \DataTypeTok{int}\NormalTok{ next}\OperatorTok{;}
\OperatorTok{\};}
\KeywordTok{typedef} \KeywordTok{struct}\NormalTok{ Node SLinkList}\OperatorTok{[}\NormalTok{MaxSize}\OperatorTok{];}
\end{Highlighting}
\end{Shaded}

即：\uln{可用 \texttt{SLinkList} 定义“一个长度为 \texttt{MaxSize} 的
\texttt{Node} 型数组”}。\texttt{SLinkList\ a;} 和
\texttt{struct\ Node\ a{[}MaxSize{]};} 等价。

\slidefig{assets/2.3.5-p05.jpg}{0.65}{课件：两种定义方式等价}

\textbf{为什么课本这么写}：和 2.3.1 的 \texttt{LinkList} /
\texttt{LNode\ *} 完全相通。用前一种写法声明，从代码上看 \texttt{a}
是一个 \texttt{Node} 型数组；用 \texttt{SLinkList\ a}
声明，效果一样，但一看就知道 \uln{\texttt{a} 是一个静态链表}。

\hypertarget{ux8001ux5e08ux7684ux9a8cux8bc1ux7a0bux5e8f}{%
\subsection{老师的验证程序}\label{ux8001ux5e08ux7684ux9a8cux8bc1ux7a0bux5e8f}}

\begin{Shaded}
\begin{Highlighting}[]
\PreprocessorTok{\#define MaxSize }\DecValTok{10}
\KeywordTok{struct}\NormalTok{ Node }\OperatorTok{\{}
    \DataTypeTok{int}\NormalTok{ data}\OperatorTok{;}
    \DataTypeTok{int}\NormalTok{ next}\OperatorTok{;}
\OperatorTok{\};}
\KeywordTok{typedef} \KeywordTok{struct} \OperatorTok{\{}
    \DataTypeTok{int}\NormalTok{ data}\OperatorTok{;}
    \DataTypeTok{int}\NormalTok{ next}\OperatorTok{;}
\OperatorTok{\}}\NormalTok{ SLinkList}\OperatorTok{[}\NormalTok{MaxSize}\OperatorTok{];}

\DataTypeTok{void}\NormalTok{ testSLinkList}\OperatorTok{()} \OperatorTok{\{}
    \KeywordTok{struct}\NormalTok{ Node x}\OperatorTok{;}
\NormalTok{    printf}\OperatorTok{(}\StringTok{"sizeX=}\SpecialCharTok{\%d\textbackslash{}n}\StringTok{"}\OperatorTok{,} \KeywordTok{sizeof}\OperatorTok{(}\NormalTok{x}\OperatorTok{));}   \CommentTok{// 8}
    \KeywordTok{struct}\NormalTok{ Node a}\OperatorTok{[}\NormalTok{MaxSize}\OperatorTok{];}
\NormalTok{    printf}\OperatorTok{(}\StringTok{"sizeA=}\SpecialCharTok{\%d\textbackslash{}n}\StringTok{"}\OperatorTok{,} \KeywordTok{sizeof}\OperatorTok{(}\NormalTok{a}\OperatorTok{));}   \CommentTok{// 80}
\NormalTok{    SLinkList b}\OperatorTok{;}
\NormalTok{    printf}\OperatorTok{(}\StringTok{"sizeB=}\SpecialCharTok{\%d\textbackslash{}n}\StringTok{"}\OperatorTok{,} \KeywordTok{sizeof}\OperatorTok{(}\NormalTok{b}\OperatorTok{));}   \CommentTok{// 80}
\OperatorTok{\}}
\end{Highlighting}
\end{Shaded}

运行结果
\texttt{sizeX=8}、\texttt{sizeA=80}、\texttt{sizeB=80}：一个结点含两个
4B 的 \texttt{int}，共 8B；\texttt{b} 和 \texttt{a} 一样大，说明
\texttt{SLinkList\ b}
\uln{相当于定义了一个长度为 \texttt{MaxSize} 的 \texttt{Node} 型数组}。

\begin{bookbox}

王道书第 38
页：静态链表是用\textbf{数组}来模拟线性表的链式存储结构。每个结点包含
\texttt{data} 域和 \texttt{next}
域；与动态链表不同，这里的指针实际上是结点在数组中的相对地址（数组下标），也称\textbf{游标}。类似于顺序表，静态链表也需要预先分配一块连续的内存空间。书上的结构类型定义和课件相同（书上
\texttt{MaxSize} 为 50）。

\end{bookbox}

\slidefig{assets/2.3.5-book-fig2.14.png}{0.75}{王道书图 2.14：(a) 静态链表示例；(b) 静态链表对应的单链表}

\hypertarget{ux4e09ux7b80ux8ff0ux57faux672cux64cdux4f5cux7684ux5b9eux73b0}{%
\section{三、简述基本操作的实现}\label{ux4e09ux7b80ux8ff0ux57faux672cux64cdux4f5cux7684ux5b9eux73b0}}

\tocback

\Needspace{4\baselineskip}

\cuenote{怎么初始化？查找是 \(O(1)\)
吗？插入、删除分几步？怎么判断一个结点是空闲的？}老师只讲思路，感兴趣的可以自己动手实现代码。

\hypertarget{ux521dux59cbux5316}{%
\subsection{初始化}\label{ux521dux59cbux5316}}

单链表初始化时把头结点的 \texttt{next} 指向
\texttt{NULL}；对应到静态链表，\uln{把 \texttt{a{[}0{]}} 的 \texttt{next} 设为 \(-1\)}（\(-1\)
相当于 \texttt{NULL}，不指向任何元素）。

\hypertarget{ux67e5ux627e}{%
\subsection{查找}\label{ux67e5ux627e}}

只能从头结点出发，顺着游标记录的线索依次往后找，直到找到想要的结点。所以找某个\textbf{位序}的结点，时间复杂度
\hlred{\(O(n)\)}。

\begin{pitfallbox}

注意说的是“某个\textbf{位序}的结点”，而不是“某个\textbf{数组下标}的结点”。位序是各结点在\textbf{逻辑上}的顺序；数组下标只反映各结点在\textbf{物理上}的顺序。

\end{pitfallbox}

\hypertarget{ux63d2ux5165ux4f4dux5e8fux4e3a-i-ux7684ux7ed3ux70b9}{%
\subsection{插入位序为 i
的结点}\label{ux63d2ux5165ux4f4dux5e8fux4e3a-i-ux7684ux7ed3ux70b9}}

\begin{enumerate}
\def\labelenumi{\arabic{enumi}.}
\tightlist
\item
  找到一个空的结点，存入数据元素（可以按某种算法，比如从上到下扫描）；
\item
  从头结点出发，找到位序为 \(i-1\) 的结点；
\item
  修改新结点的 \texttt{next}；
\item
  修改 \(i-1\) 号结点的 \texttt{next}。
\end{enumerate}

老师的例子：表中已有 4 个元素，要在位序 5（表尾）插入新元素
\(e_5\)，假设存到下标 4：

{\rowcolors{2}{wdblue!5}{white}\arrayrulecolor{wdblue!35}
\begin{xltabular}{\linewidth}{>{\hsize=0.600\hsize\raggedright\arraybackslash}X>{\hsize=1.200\hsize\raggedright\arraybackslash}X>{\hsize=1.200\hsize\raggedright\arraybackslash}X}
\rowcolor{wdblue}\textcolor{white}{\bfseries 下标} & \textcolor{white}{\bfseries 插入前 \texttt{(data,\ next)}} & \textcolor{white}{\bfseries 插入后 \texttt{(data,\ next)}} \\
\endfirsthead
0 & 头, 2 & 头, 2 \\ \hline
1 & \(e_2\), 6 & \(e_2\), 6 \\ \hline
2 & \(e_1\), 1 & \(e_1\), 1 \\ \hline
3 & \(e_4\), \(-1\) & \(e_4\), \hlred{\textbf{4}} \\ \hline
4 & 空 & \hlred{\(e_5\), \(-1\)} \\ \hline
6 & \(e_3\), 3 & \(e_3\), 3 \\ \hline
\end{xltabular}\par\vspace{4pt}}

前驱是第 4 个元素 \(e_4\)（下标 3），把它的 \texttt{next} 改成
4，指向新结点的存放位置；新结点成为表尾，\texttt{next} 设为 \(-1\)。

\hypertarget{ux600eux4e48ux5224ux65adux7ed3ux70b9ux662fux7a7aux95f2ux7684}{%
\subsection{怎么判断结点是空闲的}\label{ux600eux4e48ux5224ux65adux7ed3ux70b9ux662fux7a7aux95f2ux7684}}

刚才找空结点时，我们是用肉眼看出 4、5、7、8、9
没存数据。但从计算机的视角，内存中任何地方都有数据，只不过是\textbf{脏数据}。所以要让计算机识别空闲结点：

\begin{itemize}
\tightlist
\item
  \textbf{初始化时}还要把空闲结点的 \texttt{next} 设为某个特殊值，比如
  \hlred{\(-2\)}，写代码时判断 \texttt{next\ ==\ -2}
  就知道它是空闲的（这是前面初始化漏掉的一步）；
\item
  \textbf{删除时}除了修改游标，也要把被删除结点的 \texttt{next} 设为
  \(-2\)，表示它已经被回收。
\end{itemize}

\hypertarget{ux5220ux9664ux67d0ux4e2aux7ed3ux70b9}{%
\subsection{删除某个结点}\label{ux5220ux9664ux67d0ux4e2aux7ed3ux70b9}}

\begin{enumerate}
\def\labelenumi{\arabic{enumi}.}
\tightlist
\item
  从头结点出发找到前驱结点；
\item
  修改前驱结点的游标；
\item
  被删除结点的 \texttt{next} 设为 \(-2\)。
\end{enumerate}

\slidefig{assets/2.3.5-p08.jpg}{0.65}{课件：简述基本操作的实现}

\begin{bookbox}

王道书第 38 页：静态链表以 \texttt{next\ ==\ -1}
作为其结束标志。插入、删除操作与动态链表类似，只需修改指针（数组下标），而无须移动元素。尽管静态链表在灵活性上不及动态链表，但在一些不支持指针的编程语言（如
Basic）中，它提供了一种巧妙的设计方案。

\end{bookbox}

\newpage

\hypertarget{ux56dbux77e5ux8bc6ux56deux987e}{%
\section{四、知识回顾}\label{ux56dbux77e5ux8bc6ux56deux987e}}

\tocback

\slidefig{assets/2.3.5-p09.jpg}{0.65}{课件：知识回顾与重要考点}

\begin{summarybox}

\textbf{静态链表}：用数组的方式实现的链表。存储空间是一整片连续空间，但在这片空间内，逻辑上相邻的元素物理上也可以不相邻，先后关系靠\textbf{游标}（下一个结点的数组下标）表示。0
号结点充当头结点；游标 \(-1\) 表示表尾；可以用 \(-2\)
等特殊值标记空闲结点。

\begin{itemize}
\tightlist
\item
  \textbf{优点}：增、删操作不需要大量移动元素，只需修改相关结点的游标；
\item
  \textbf{缺点}：不能随机存取，只能从头结点开始依次往后查找；\hlred{容量固定不可变}，声明后最大容量就定死了，不可扩展；
\item
  \textbf{适用场景}：①
  不支持指针的低级语言（早期用它实现和单链表同样的功能）；②
  数据元素数量固定不变的场景，如操作系统文件管理中的\textbf{文件分配表
  FAT}，本质上就是一个静态链表（老师建议先在操作系统课本上做个记号，学到文件管理时回头体会）。
\end{itemize}

\texttt{typedef\ struct\ \{…\}\ SLinkList{[}MaxSize{]};} 让
\texttt{SLinkList} 代表“长度为 \texttt{MaxSize}
的结点数组”，\texttt{SLinkList\ a;} 等价于
\texttt{struct\ Node\ a{[}MaxSize{]};}，写法上强调 \texttt{a}
是一个静态链表。

\textbf{自测}：① 静态链表的游标和单链表的指针有什么异同？② 结点
8B、起始地址 \texttt{addr}，下标 5 的结点地址是多少？③
\texttt{SLinkList\ a;} 等价于什么声明？④ 为什么初始化时要把空闲结点的
\texttt{next} 设为 \(-2\)？⑤ 静态链表找位序为 \(i\)
的结点，时间复杂度是多少？

\end{summarybox}
\end{document}
